\chapter{Recent Methods}

\section{Linear DiD with staggered timing}

\texttt{lpdid} (Local Projections DiD) and \texttt{eventstudy} estimate treatment
effects in designs with staggered adoption. \texttt{eventstudy} constructs
event-time dummies relative to treatment and drops the reference period.

\begin{lstlisting}[caption={Event study}]
let m_e = eventstudy(y ~ x, df, time=year, ref=-1,
                     min=-5, max=5)
print(m_e)
\end{lstlisting}

The coefficients on each event-time dummy show the average outcome path
relative to the treatment period.

\section{Synthetic DiD and CUPED}

Synthetic DiD combines the synthetic-control idea with difference-in-
differences. \texttt{synthdid} estimates unit and time weights from the pre-
treatment outcomes and reports the ATT.

\begin{lstlisting}[caption={Synthetic DiD}]
let m_sd = synthdid(df, y, treated, 10, unit="state", period="year")
print(m_sd)
\end{lstlisting}

CUPED adjusts the outcome with pre-treatment covariates to reduce
variance in randomised experiments.

\begin{lstlisting}[caption={CUPED}]
let m_c = cuped(y ~ x, df, treated="treated")
print(m_c)
\end{lstlisting}

\section{MIDAS regression}

Mixed Data Sampling (Ghysels, Santa-Clara and Valkanov, 2004) regresses
a low-frequency variable on high-frequency lags, with weights constrained
by an Almon polynomial.

\begin{lstlisting}[caption={MIDAS}]
let m_midas = midas(y ~ x, df, freq=3, lags=12, poly=2)
print(m_midas)
\end{lstlisting}

It is useful for nowcasting: using monthly predictors to forecast a
quarterly target.
